\documentclass{article}

\usepackage{amsmath}
\usepackage{amsfonts}
\usepackage{geometry}
\geometry{margin=1in}

\title{Technical Summary: Parallel and Distributed Computing}
\author{Documentation Synthesis}
\date{\today}

\begin{document}

\maketitle

\section{Introduction}
The field of Parallel and Distributed Computing (PDC) encompasses the study and application of multiple simultaneous computing resources to solve a task. This document outlines the fundamental distinctions between concurrency and parallelism, explores the execution model of multi-threading, and identifies common synchronization challenges and solutions.

\section{Concurrency vs. Parallelism}
A fundamental distinction exists between how tasks are managed and how they are executed.

\subsection{Concurrency}
Concurrency is the property of a system where multiple tasks are in progress at the same time. It is a structural approach to software design, allowing the system to deal with multiple things at once. In a single-core environment, this is achieved through context switching, where the CPU interleaves the execution of different tasks.

\subsection{Parallelism}
Parallelism is the property where multiple tasks are physically executed at the exact same moment. This requires hardware with multiple processing units, such as multi-core CPUs or GPUs. Parallelism is a subset of concurrency focused on increasing execution speed.

\subsection{Comparative Analysis}
\begin{center}
  \begin{tabular}{|l|p{5cm}|p{5cm}|}
    \hline
    \textbf{Feature} & \textbf{Concurrency} & \textbf{Parallelism} \\
    \hline
    Focus & Task management and structure. & Task execution and performance. \\
    \hline
    Hardware & Possible on a single-core CPU. & Requires multi-core or multi-processor hardware. \\
    \hline
    Mechanism & Interleaving/Context switching. & Simultaneous execution. \\
    \hline
  \end{tabular}
\end{center}

\section{The Multi-threading Model}
Multi-threading is an execution model that allows a single process to contain multiple sequences of instructions (threads) running in the same address space.

\begin{itemize}
  \item \textbf{Process:} An instance of a computer program that is being executed, containing its own memory space.
  \item \textbf{Thread:} The smallest unit of processing that can be performed in an OS. Threads within the same process share the heap but have their own stack and registers.
  \item \textbf{Context Switch:} The process of storing the state of a process or thread so that it can be restored and execution resumed from the same point later.
\end{itemize}

\section{Challenges in Concurrent Systems}
Writing software that manages shared state across multiple threads introduces several critical failure modes.

\subsection{Race Conditions}
A race condition occurs when the output of a program depends on the relative timing of threads. If two threads access and modify the same variable concurrently without proper synchronization, the final value may be inconsistent.

\subsection{Deadlock}
A deadlock is a state where a set of threads are blocked because each thread is waiting for a resource held by another thread in the set. The four necessary conditions for deadlock are:
\begin{enumerate}
  \item Mutual Exclusion
  \item Hold and Wait
  \item No Preemption
  \item Circular Wait
\end{enumerate}

\section{Synchronization Primitives}
To mitigate the risks of concurrent execution, various synchronization mechanisms are employed.

\subsection{Mutex (Mutual Exclusion)}
A mutex is a locking mechanism used to synchronize access to a resource. Only one thread can acquire the mutex at a time. Other threads attempting to acquire it will block until it is released.

\subsection{Semaphores}
A semaphore is a signaling mechanism that uses a counter to control access to a shared resource.
\begin{itemize}
  \item \textbf{Binary Semaphore:} Operates similarly to a mutex (values 0 or 1).
  \item \textbf{Counting Semaphore:} Allows a set number of threads to access a resource pool.
\end{itemize}

\subsection{Atomic Operations}
Atomic operations are instructions that complete in a single step relative to other threads. They are uninterruptible, ensuring that no thread sees the operation in a partially completed state.

\section{Conclusion}
While concurrency provides a way to structure programs to handle multiple tasks efficiently (particularly I/O-bound tasks), parallelism provides the means to accelerate computation-heavy (CPU-bound) tasks. Understanding the interplay between these concepts, the underlying threading models, and the necessary synchronization primitives is essential for systems-level development.

\end{document}